\centerline{\bf \Huge Можно ли? I}

\begin{flushright}
        \scriptsize{Можно все, что хочешь, когда хочешь только то, что следует.}
\end{flushright}

\problem{\textbf{1}} Может ли в одном месяце быть 5 понедельников? А может 6 ?

\problem{\textbf{2}} Можно ли в прямоугольную таблицу поставить числа так, чтобы в каждом столбце сумма была положительна, а в каждой строке - отрицательна?

\problem{\textbf{3}} Можно ли подобрать такие два последовательных натуральных числа, что сумма цифр каждого из них делится на 4?

\problem{\textbf{4}} Может ли и сумма, и произведение нескольких натуральных чисел быть равными 111? A 101?

\begin{flushright} \textbf{2 балла}\\ 
\end{flushright}

\problem{\textbf{5}} Можно ли выбить 100 очков несколькими выстрелами по мишени, в которой есть зоны $9,12,15,18,24$ и 47 очков?

\problem{\textbf{6}} Можно ли внутри выпуклого пятиугольника отметить 18 точек так, чтобы внутри каждого из десяти треугольников, образованных его вершинами, отмеченных точек было поровну?

\problem{\textbf{7}} Расул Владимирович сложил на столе из 9 квадратов и 19 равносторонних треугольников (не накладывая их друг на друга) многоугольник. Мог ли периметр
этого многоугольника оказаться равным 15 см, если стороны всех квадратов и треугольников
равны 1 см?


\begin{flushright} \textbf{4 балла}\\ 
\end{flushright}
